% titel: Blatt 3, A2: Masse am Faden und Rohr -- Drehimpuls
% skript: Kap. Drehbewegung, Drehimpulserhaltung

% Drehimpulserhaltung, Zentripetalkraft, Energie bei Radiusänderung.

Zieht man die Punktmasse durch das Rohr nach innen, bleibt der
Drehimpuls $L$ erhalten, weil keine aeussere Kraft ein Drehmoment
bezüglich der Achse ausübt.

\begin{werkzeugkasten}[Werkzeuge]
\textbf{Drehimpuls} einer Punktmasse auf Kreisbahn:
\begin{equation}
  L=m\,v\,r=m\,r^2\,\omega\;.
\end{equation}

\textbf{Drehimpulserhaltung} (kein aeusseres Drehmoment):
\begin{equation}
  m\,r_1^2\,\omega_1=m\,r_2^2\,\omega_2
  \;\Rightarrow\;\omega_2=\omega_1\left(\frac{r_1}{r_2}\right)^2\;.
\end{equation}

\textbf{Zentripetalkraft} (Fadenspannung):
\begin{equation}
  F=m\,\frac{v^2}{r}=m\,r\,\omega^2\;.
\end{equation}

\textbf{Kinetische Energie}:
\begin{equation}
  E_{\rm kin}=\tfrac12 m v^2=\tfrac12 m r^2\omega^2.
\end{equation}
\textbf{Arbeitssatz:} die von der aeu{\ss}eren Zugkraft verrichtete Arbeit
gleicht der Aenderung der kinetischen Energie, $W=\Delta E_{\rm kin}$.
\end{werkzeugkasten}

\begin{werkzeugkasten}[Vorgehen]
\begin{itemize}
  \item Für $\omega_2$ nur die Drehimpulserhaltung nutzen.
  \item Die Kreisfrequenz geht quadratisch ein: kleiner Radius liefert
        stark verstärktes $\omega_2$.
  \item Arbeit = Energiebilanze ($\tfrac12 m r^2\omega^2$) vorher/nachher.
\end{itemize}
\end{werkzeugkasten}

\begin{achtung}[Typische Fallen]
\begin{itemize}
  \item Drehimpulserhaltung sagt nichts über die Energie: die kinetische
        Energie steigt, weil die Zugkraft Arbeit verrichtet.
  \item $\omega\propto 1/r^2$ (nicht $1/r$); $\omega^2$ noch steiler.
  \item Fadenspannung wächst mit $\omega^2 r$; kleine Radien -> grosse
        Spannung.
\end{itemize}
\end{achtung}
